\section{Apparatus and Methods}
\label{sec:apparatus}

\paragraph{Models.}
All games are played by two LLM agents drawn from four open-weight models served through
hosted inference (LiteLLM~$\to$~OpenRouter): \textbf{Qwen2.5-72B} and \textbf{DeepSeek-V3}
(grouped under the label ``CN-origin''), and \textbf{Llama-3-70B} and
\textbf{Mistral-Large} (``Western-origin''), plus \textbf{GPT-4o} as a frontier anchor. We
emphasize at the outset that ``origin'' is a confounded proxy bundling family, size,
quantization, serving stack, fine-tuning, and training language; we treat it accordingly
(\S\ref{sec:origin}). Both spine models passed an \textbf{A5 validity gate} (sane responses
to classical strategies such as tit-for-tat and always-defect) before being admitted to the
confirmatory spine. Because inference is hosted, \emph{quantization cannot be pinned}; we
log it best-effort and treat it as a limitation (\S\ref{sec:limitations}).

\paragraph{The communication / affordance ladder.}
The core manipulation is a four-rung ladder that varies the two affordances:
\begin{description}[leftmargin=2.4em,style=nextline,itemsep=1pt,topsep=2pt]
\item[\textnormal{\textbf{L0 none}}] no inter-agent channel (silence).
\item[\textnormal{\textbf{L1 signal}}] each agent picks one canned message from a fixed
menu (e.g.\ ``I will cooperate'' / ``I will defect''). High accountability, near-zero
semantic bandwidth.
\item[\textnormal{\textbf{L2 observed}}] free-text messages that the principal / PI
\emph{can} read (high bandwidth, observable).
\item[\textnormal{\textbf{L3 private}}] free-text messages on a back-channel the principal
\emph{cannot} read (high bandwidth, unobservable).
\end{description}
\noindent L0$\to$L1$\to$L2 increases semantic bandwidth at fixed observability; L2$\to$L3
holds bandwidth fixed and removes observability. This factoring lets us attribute effects to
\emph{content} versus \emph{accountability} rather than to ``communication'' in the
aggregate. A companion content/authorship ablation (C1) further separates self-authorship
from message content.

\paragraph{Primary outcome: lock-in proportion with Wilson intervals.}
Matches are strongly bimodal---a pair either settles into sustained mutual cooperation or
into sustained defection---so a round-mean misrepresents the central tendency. We therefore
define the primary cooperation outcome as the \textbf{match-level lock-in proportion}: the
fraction of matches whose mean cooperation exceeds $C>0.8$. We report it with a
\textbf{Wilson 95\% score interval} \citep{wilson1927probable}, which is well-behaved for
proportions near $0$ and $1$ and at small $n$. \textbf{Refusals are coded as outcomes,
never silently dropped}; API/transport errors are excluded from behavioral metrics and
tracked separately as an \texttt{api\_error\_rate}. Null claims are ``earned'' with two
one-sided equivalence tests (TOST) \citep{lakens2017tost,lakens2018equivalence} rather than
read off non-significance.

\paragraph{The collusion index $K$, and why $K>1$ is real.}
For market games we summarize the realized price level with a collusion index
\[
K = \frac{p - p_{\text{competitive}}}{p_{\text{monopoly}} - p_{\text{competitive}}},
\]
normalized so that $K{=}0$ is the static competitive (Bertrand/Nash) benchmark and
$K{=}1$ is the static joint-monopoly benchmark. Crucially, \textbf{$K>1$ is expected and we
do not clip it.} In finite-horizon, discretized repeated play, agents can sustain prices
\emph{above} the one-shot monopoly point using \emph{dynamic punishment} strategies (the
threat of reverting to a price war disciplines deviation), exactly as the folk theorem for
repeated games predicts \citep{friedman1971noncooperative,fudenberg1986folk}. We therefore report values such as $K{=}1.23$ as a genuine
\textbf{supra-monopoly} regime rather than as an artifact or a normalization bug. Clipping
$K$ at $1$ would erase the most principal-harming behavior we observe.

\paragraph{Anti-memorization: novel payoffs and novel demand.}
A standing worry for LLM game-play is that models recall canonical games and benchmark
results (e.g.\ Calvano/Fish pricing) rather than reason. We therefore run every spine
experiment on both \textbf{canonical} and \textbf{novel} payoff/demand structures---a
demand curve and payoff matrix constructed to be absent from any plausible training
corpus. Effects that survive the novel condition are attributed to in-context reasoning,
not recall; and we use the canonical-vs-novel gap to flag that magnitudes measured on
canonical structures are \emph{inflated} (\S\ref{sec:results}).

\paragraph{Two-speed protocol.}
To avoid exploratory-to-confirmatory leakage, the pre-registered spine (A1 cooperation, B1
pricing, A3 origin) is locked in \texttt{prereg/} and run on \emph{fresh held-out seeds}
(seed offset 100) only after the analysis plan is fixed; the wider $\sim$30-family catalog
is exploratory and reported as such. Every number in this paper is tagged
confirmatory~\tC{} or exploratory~\tE{}.

\paragraph{Sandbox containment for alignment scenarios.}
The alignment wave (\S\ref{sec:alignment}) places agents in a tool-using ``Lab'' (rival
research labs, a shared overseer, file/message tools). \textbf{All tools are mocked and the
Lab is an in-memory sandbox}; a runtime \texttt{assert\_contained()} check enforces that no
action escapes the process (no network egress, no real files, no real recipients).
Misaligned acts are instrumented against ground truth (e.g.\ whether confidential findings
actually crossed to the rival), so we measure \emph{functional} harm, not surface wording.
